\section{模型评价与结论}

\subsection{模型优点}
\begin{enumerate}
  \item 建模链条与四个问题一一对应，避免使用与题目目标无关的评价、优化或时间序列模型。
  \item 同时考虑无效数据、未检出值、重复采样和成分闭合效应，数据处理较为完整。
  \item 分类模型优先选择低复杂度、可解释方法，并以按文物分组的交叉验证降低信息泄漏。
  \item 对未知样品和亚类划分均进行扰动敏感性分析，不只给出单一确定结论。
\end{enumerate}

\subsection{模型局限}
\begin{enumerate}
  \item 缺少充足的同一文物风化点—未风化点配对数据，风化前成分恢复主要依赖横截面规律，因果解释有限。
  \item 埋藏时间、环境酸碱度和温湿度等变量未被观测，可能形成遗漏变量偏差。
  \item 部分微量成分未检出比例较高，其相关关系和聚类贡献存在不确定性。
\end{enumerate}

\subsection{主要结论}
\begin{enumerate}
  \item \todo{问题一结论：风化与哪些属性显著相关，关键成分如何变化，恢复模型误差如何。}
  \item \todo{问题二结论：高钾/铅钡的核心分类规律，各自划分为哪些亚类。}
  \item \todo{问题三结论：A1--A8的类型及边界样品。}
  \item \todo{问题四结论：两类玻璃的主要成分关联及差异。}
\end{enumerate}
